\documentclass[12pt,a4paper]{article}
% Linh hoạt theo engine: pdflatex (cần vntex/T5) hoặc xelatex/lualatex (Unicode trực tiếp).
\usepackage{iftex}
\ifPDFTeX
  \usepackage[utf8]{inputenc}
  \IfFileExists{t5enc.def}{\usepackage[T5]{fontenc}}{\usepackage[T1]{fontenc}}
\else
  \usepackage{fontspec}
\fi
% Chỉ nạp babel tiếng Việt nếu có sẵn (vntex/babel-vietnamese); nếu không, fontspec/Unicode tự render.
\IfFileExists{vietnamese.ldf}{\usepackage[vietnamese]{babel}}{%
  \renewcommand{\contentsname}{Mục lục}%
  \renewcommand{\refname}{Tài liệu tham khảo}%
  \renewcommand{\tablename}{Bảng}%
  \renewcommand{\figurename}{Hình}%
}
\usepackage{amsmath, amssymb, amsthm, mathtools}
\usepackage{booktabs}
\usepackage{graphicx}
\usepackage[hidelinks]{hyperref}
\usepackage{geometry}
\usepackage{setspace}
\usepackage{enumitem}
\usepackage{array}
\usepackage{longtable}
\usepackage{multirow}
\usepackage{cite}
\usepackage{microtype}
\usepackage{titlesec}
\usepackage{float}
\usepackage{listings}
\usepackage{xcolor}

\geometry{margin=1in}
\setstretch{1.25}
\hypersetup{
    colorlinks=true,
    linkcolor=black,
    urlcolor=cyan,
    citecolor=red
}

\lstset{
    basicstyle=\ttfamily\footnotesize,
    keywordstyle=\color{blue},
    commentstyle=\color{gray},
    stringstyle=\color{teal},
    breaklines=true,
    frame=single,
    showstringspaces=false,
    language=Python,
    columns=fullflexible,
}

% =========================================================
% Theorem environments
% =========================================================
\newtheorem{theorem}{Định lý}[section]
\newtheorem{lemma}{Bổ đề}[section]
\newtheorem{proposition}{Mệnh đề}[section]
\newtheorem{corollary}{Hệ quả}[section]
\newtheorem{definition}{Định nghĩa}[section]
\newtheorem{remark}{Nhận xét}[section]

% =========================================================
% Title
% =========================================================
\title{
    \vspace{-1.8cm}
    \begin{center}
        \vspace{0.5cm}
        \LARGE \textbf{U-MobileViT cho phân vùng thời gian thực trên thiết bị biên} \\
        \LARGE \textbf{và tính mới INLA: Inverted Nonlinear Low-rank Attention} \\
        \Large \textit{Báo cáo dự án: kiến trúc, tối ưu hóa edge, phân tích cơ chế kiểm soát rank collapse} \\
        \vspace{0.5cm}
    \end{center}
}

\author{Nguyễn Ngọc Bình An, Hoàng Thị Linh Hương}
\date{\today}

\begin{document}

\maketitle

\begin{abstract}
Báo cáo này trình bày dự án \textbf{U-MobileViT} --- một mạng phân vùng (segmentation) gọn nhẹ
theo phong cách U-Net lai MobileViT, được thiết kế để chạy thời gian thực trên thiết bị biên
(Jetson, ARM). Phần thứ nhất mô tả kiến trúc, các tối ưu hóa cấp hệ thống đã thực hiện
(chia sẻ head đa nhiệm, lượng tử hóa, xuất ONNX/TensorRT) cùng một thư viện khối thị giác và
transformer tái sử dụng. Phần thứ hai đề xuất \textbf{tính mới} của dự án:
\emph{Inverted Nonlinear Low-rank Attention} (INLA), một cơ chế attention nhẹ áp dụng phép nâng
đặc trưng phi tuyến theo tinh thần inverted bottleneck \emph{trước} bước tổng hợp linear attention,
nhằm làm chậm hiện tượng suy giảm hạng (rank collapse) và thoái hóa phổ trong các backbone nhỏ
trong khi vẫn giữ độ phức tạp tuyến tính theo số token.

Nghiên cứu theo đuổi ba lớp kết quả: (i) mô tả cơ chế rank collapse / spectral degeneration;
(ii) xây dựng các mệnh đề và phân tích cục bộ về tác động của feature lifting lên effective rank,
singular value decay và conditioning; (iii) hiện thực INLA như một module thật trong codebase và
kiểm chứng sơ bộ bằng thực nghiệm. Chúng tôi báo cáo các số liệu \emph{đo thật} trên mã nguồn của
dự án: độ phức tạp được xác nhận tuyến tính theo số token, chi phí tham số tăng thêm khiêm tốn,
công cụ chẩn đoán phổ được kiểm định đúng đắn, và một thí nghiệm thăm dò rank collapse minh họa cả
trường hợp cơ chế lộ rõ (attention thuần) lẫn trường hợp bị che khuất (có residual). Tinh thần
xuyên suốt là trung thực về điều kiện áp dụng: chỉ ra khi nào nonlinear lifting thực sự giúp ích và
khi nào nó chỉ tạo overhead.
\end{abstract}

\tableofcontents
\newpage

% =========================================================
\section{Mở đầu}

\subsection{Bối cảnh dự án}
Dự án xuất phát từ nhu cầu thực tế: triển khai một mô hình phân vùng ảnh chất lượng tốt nhưng
\emph{đủ nhẹ để chạy thời gian thực trên thiết bị biên} (Jetson Nano/Orin, ARM, NPU). Transformer
và attention đã trở thành nền tảng quan trọng của học sâu nhờ khả năng mô hình hóa phụ thuộc xa và
tương tác toàn cục giữa các token \cite{vaswani2017attention}. Tuy nhiên self-attention tiêu chuẩn
có độ phức tạp và bộ nhớ bậc hai theo số token, trở thành nút thắt khi triển khai trên phần cứng
hạn chế tài nguyên.

Nhiều hướng \emph{efficient attention} và \emph{linear attention} đã được đề xuất nhằm xấp xỉ hoặc
thay thế self-attention bằng các phép tính chi phí thấp hơn \cite{katharopoulos2020transformers,
shen2020efficient, choromanski2021performer, wang2020linformer, xiong2021nystromformer,
qin2022cosformer}. Điểm chung là tránh xây dựng trực tiếp ma trận tương tác token--token kích thước
$N\times N$, dựa vào kernel feature maps, low-rank approximation, random features hoặc các phép phân
rã khác để giảm chi phí xuống gần tuyến tính. Backbone của dự án (U-MobileViT) chính là một hiện
thực theo hướng này: kết hợp convolution nhẹ (MobileNetV2) với linear attention kiểu MobileViTv2
\cite{mobilenetv2, mobilevitv1, mobilevitv2}.

\subsection{Vấn đề cốt lõi}
Khi chi phí được giảm mạnh, câu hỏi quan trọng không chỉ là ``có chạy nhanh hơn không'' mà còn là
``chất lượng biểu diễn còn lại như thế nào''. Trọng tâm của phần nghiên cứu là hiện tượng
\textbf{rank collapse} và \textbf{spectral degeneration} trong linear attention. Khi ánh xạ feature
map bị hạn chế, nhiều query khác nhau có thể bị chiếu vào các hướng gần giống nhau; lặp lại qua
nhiều lớp, attention map kém đa dạng, singular values suy giảm nhanh, và biểu diễn token đồng nhất
hóa dần. Hiện tượng này đặc biệt đáng lo trong backbone nhẹ, nơi chiều ẩn, số head và ngân sách
tính toán đều bị giới hạn.

\subsection{Câu hỏi nghiên cứu trung tâm}
\begin{quote}
\textbf{Liệu một cơ chế inverted nonlinear feature lifting có thể làm chậm tốc độ suy giảm hạng và
ổn định cấu trúc phổ của linear attention theo chiều sâu mạng, trong khi vẫn duy trì độ phức tạp
tuyến tính và tính hiệu quả cho mô hình gọn nhẹ?}
\end{quote}

\subsection{Định vị học thuật}
Báo cáo này không phải pure theory, cũng không phải thuần ``thay module rồi báo cáo accuracy''.
Định vị phù hợp là \textit{mechanism-level research with theoretical analysis and empirical
validation}: cơ chế là câu hỏi chính, chứng minh là công cụ làm rõ cơ chế, thực nghiệm là bước
kiểm tra tính đúng đắn và hữu ích.

\subsection{Đóng góp}
\begin{itemize}
    \item \textbf{Hệ thống:} một backbone U-MobileViT chạy được thời gian thực, kèm tối ưu hóa edge
    (chia sẻ head đa nhiệm giảm 26\% FLOPs, xuất ONNX/TorchScript, lượng tử hóa) --- xem
    Mục~\ref{sec:project}.
    \item \textbf{Kỹ thuật phần mềm:} một thư viện khối thị giác \& transformer tái sử dụng
    (\texttt{cv\_nets.blocks}) và công cụ chẩn đoán phổ (\texttt{cv\_nets.utils.spectral}).
    \item \textbf{Cơ chế \& lý thuyết:} làm rõ rank collapse trong linear attention và xây dựng các
    mệnh đề/định lý cục bộ về tác động của nonlinear lifting --- Mục~\ref{sec:theory}.
    \item \textbf{Tính mới INLA:} hiện thực INLA như một module thật, tích hợp vào codebase, kèm
    kiểm chứng thực nghiệm sơ bộ --- Mục~\ref{sec:method},~\ref{sec:impl},~\ref{sec:exp}.
\end{itemize}

% =========================================================
\section{Dự án U-MobileViT: kiến trúc và tối ưu hóa edge}
\label{sec:project}

\subsection{Tổng quan kiến trúc}
U-MobileViT là một mạng encoder--decoder dạng U-Net cho phân vùng dày đặc (dense prediction).
Encoder gồm một \emph{stem} ba lớp convolution stride-2 (hạ mẫu $\times 8$) rồi ba \emph{stage},
mỗi stage có một lớp hạ mẫu depthwise stride-2 và các lớp MobileViT (khối cục bộ depthwise-separable
+ khối toàn cục transformer dùng \emph{separable/linear attention}). Decoder gồm các stage nâng mẫu
với transformer cross-attention nối tắt (skip) tới encoder; cuối cùng là \emph{segmentation head}
phục hồi độ phân giải và phân loại pixel. Mô hình hỗ trợ đa nhiệm (nhiều head song song) hoặc đơn
nhiệm.

\subsection{Linear/separable attention trong backbone}
Thay vì tính trực tiếp $\operatorname{Softmax}(\mathbf{Q}\mathbf{K}^T)$, backbone dùng dạng
separable attention với một \emph{context vector} toàn cục, đạt độ phức tạp tuyến tính theo số
token --- nền tảng để chèn attention vào mô hình thời gian thực. Chính đặc điểm này khiến backbone
là đối tượng lý tưởng để nghiên cứu rank collapse và đề xuất INLA.

\subsection{Tối ưu hóa cấp hệ thống đã thực hiện}
Một số tối ưu hóa đã được hiện thực và đo đạc trong dự án:
\begin{itemize}
    \item \textbf{Chia sẻ head đa nhiệm:} thay vì chạy lại toàn bộ pipeline nâng mẫu + tinh chỉnh
    cho \emph{mỗi} tác vụ, ta chạy phần này \emph{một lần} và chỉ tách nhánh ở bộ phân loại nhẹ cho
    mỗi tác vụ (chuẩn thiết kế YOLOP/HybridNets). Kết quả đo: tổng FLOPs giảm từ $0.301$ xuống
    $0.222$ GFLOPs ($-26\%$), riêng segmentation head giảm $\sim 45\%$; độ trễ CPU giảm $\sim 25\%$.
    \item \textbf{Công cụ xuất \& benchmark cho Jetson:} fuse Conv+BN, xuất \textbf{ONNX} (cho
    TensorRT) và TorchScript, benchmark FP32/FP16 GPU và INT8 dynamic CPU. Trên GPU tham chiếu, mô
    hình đạt $\sim 124$ FPS ở FP32, độ phân giải $320\times 320$.
    \item \textbf{Thư viện khối tái sử dụng:} \texttt{cv\_nets.blocks} cung cấp các khối chuẩn
    (ConvNormAct, MV2Block/InvertedResidual có SE, DepthwiseSeparableConv), attention nhẹ
    (SqueezeExcite, ECA, CBAM), transformer (PatchEmbed, MHSA, TransformerEncoderBlock), khối lai
    MobileViTBlock, và --- là tính mới của báo cáo này --- INLA.
\end{itemize}

% =========================================================
\section{Cơ sở lý thuyết và động cơ thiết kế của INLA}
\label{sec:theory}

\subsection{Triết lý bottleneck--expansion từ mô hình nhẹ}
MobileNetV2 cho thấy không gian trung gian rộng hơn giúp học biểu diễn tốt hơn, trong khi bottleneck
ở đầu vào/đầu ra giữ chi phí thấp \cite{mobilenetv2}. MobileViT cho thấy attention có thể nhúng vào
backbone nhẹ nếu cấu trúc được điều chỉnh cẩn thận \cite{mobilevitv1, mobilevitv2}. Từ đó nảy sinh
ý tưởng \emph{inverted nonlinear feature lifting}: trước khi aggregation toàn cục, ta mở rộng biểu
diễn lên một không gian cao hơn để làm phong phú basis mà attention vận hành trên đó.

\subsection{Rank collapse và spectral degeneration}
Nhiều nghiên cứu chỉ ra attention sâu có thể dẫn tới suy giảm hạng nghiêm trọng theo chiều sâu
\cite{dong2021rankcollapse}, và multi-head attention chịu một \emph{low-rank bottleneck} khi chiều
mỗi head quá nhỏ \cite{bhojanapalli2020low}. Có thể đo bằng: effective rank của biểu diễn, tốc độ
suy giảm singular values, spectral entropy, và các chỉ số oversmoothing. Trong linear attention,
mối lo này càng rõ vì không gian feature map đã bị ràng buộc trước.

\subsection{Tại sao cần nonlinear lifting}
Mục tiêu không phải phình mô hình vô điều kiện, mà tạo một vùng trung gian nơi mô hình học được
basis phong phú hơn trước khi bị nén lại: (i) thêm chiều tự do để các query không bị ép vào cùng
một cấu trúc; (ii) làm mềm sự co cụm của singular values; (iii) cung cấp mức phi tuyến đủ để thích
ứng dữ liệu thực mà không vượt giới hạn tính toán.

% =========================================================
\section{Mệnh đề lý thuyết và giả thuyết nghiên cứu}

\subsection{Các định nghĩa làm việc}
\begin{definition}[Hạng hiệu dụng]
Với ma trận $\mathbf{M}$ có singular values $\sigma_1 \geq \sigma_2 \geq \cdots \geq 0$, hạng hiệu
dụng được định nghĩa (Roy \& Vetterli) là $\mathrm{erank}(\mathbf{M}) = \exp\!\big(H(p)\big)$ với
$p_i = \sigma_i / \sum_j \sigma_j$ và $H$ là entropy Shannon --- đo số chiều thực sự đóng góp đáng
kể cho biểu diễn.
\end{definition}

\begin{definition}[Suy giảm phổ]
Biểu diễn có spectral degeneration mạnh nếu singular values giảm nhanh, dẫn tới một số ít chiều chi
phối phần lớn năng lượng của ma trận.
\end{definition}

\begin{definition}[Rank collapse]
Hiện tượng qua nhiều lớp, biểu diễn token mất dần độ đa dạng tuyến tính, dẫn đến effective rank
giảm và các token trở nên tương tự nhau hơn.
\end{definition}

\subsection{Mệnh đề định hướng}
\begin{proposition}[Giảm xác suất nhập nhằng biểu diễn]
Nếu attention được tính trên một feature map có bậc tự do lớn hơn, sinh bởi một lifting phi tuyến đủ
giàu, thì các query khác nhau có xu hướng tạo ra các phân phối attention tách biệt hơn so với feature
map quá thấp chiều. Đây là khả năng phân biệt cục bộ, không phải injective toàn cục.
\end{proposition}

\begin{proposition}[Cải thiện cận trên của hạng hiệu dụng]
Với bộ đệm Key--Value trong linear attention, hạng biểu diễn khả dụng bị giới hạn bởi số chiều của
không gian đặc trưng trung gian. Khi lifting mở rộng không gian lên $r$ chiều với $r > d_k$, cận trên
thực dụng của effective rank có thể tăng, làm chậm suy giảm phổ khi xếp chồng sâu hơn.
\end{proposition}

\begin{proposition}[Ổn định phổ dưới inverted lifting]
Nếu nonlinear lifting kết hợp với chuẩn hóa và regularization phù hợp, phổ singular value của
attention map có thể duy trì phân bố ít dốc hơn so với baseline tuyến tính, đặc biệt ở các lớp trung
gian của backbone sâu.
\end{proposition}

\subsection{Bổ đề kỹ thuật}
\begin{lemma}[Tổng hợp tuyến tính qua basis mở rộng]
Giả sử $\hat{\mathbf{K}} \in \mathbb{R}^{N \times r}$ và $\mathbf{V} \in \mathbb{R}^{N \times d_v}$.
Khi tính $\mathbf{S}=\hat{\mathbf{K}}^T\mathbf{V}$, mọi thông tin từ $\mathbf{V}$ chỉ đi qua $r$ trục
basis của không gian trung gian. Nếu $r$ tăng nhưng vẫn được regularization kiểm soát, mô hình có
thể lưu giữ nhiều mẫu tương tác hơn trước khi bị nén ở đầu ra.
\end{lemma}
\begin{proof}[Phác thảo]
$\hat{\mathbf{K}}^T\mathbf{V}$ là phép chiếu dữ liệu lên basis của $\hat{\mathbf{K}}$. Khi $r$ nhỏ,
basis nghèo và các mẫu tương tác bị đồng nhất hóa nhanh; khi $r$ lớn hơn, số hướng mà $\mathbf{V}$
được nén vào tăng lên. Vì $\mathbf{S}$ là đại diện nén của toàn chuỗi, số basis hữu ích tăng làm
tăng khả năng bảo toàn cấu trúc phổ.
\end{proof}

\begin{lemma}[Phi tuyến làm thay đổi độ cong không gian đặc trưng]
Nếu $\Phi_{\mathrm{INLA}}$ chứa ít nhất một thành phần phi tuyến trơn, ánh xạ từ token gốc sang không
gian attention không còn là affine thuần túy, làm tăng khả năng tách các cụm token vốn gần nhau.
\end{lemma}
\begin{proof}[Phác thảo]
Phép affine chỉ tạo các siêu phẳng song song, không đủ bẻ cong biên quyết định. Thành phần phi tuyến,
dù nhẹ, làm ánh xạ giàu hơn về hình học; các token gần nhau nhưng khác ở hướng phi tuyến có thể bị
đẩy sang các miền đặc trưng khác nhau, giảm nguy cơ nhiều query bị chiếu vào cùng một cấu trúc.
\end{proof}

\begin{theorem}[Chậm hóa rank collapse, cục bộ và có điều kiện]
Xét một block linear attention với feature map cố định và một block INLA với lifting phi tuyến, các
tham số lifting bị ràng buộc bởi chuẩn hóa/regularization đủ mạnh. Khi token đầu vào có tương quan
trung bình và không gian sau lifting có $r > d_k$, tồn tại một ngưỡng độ sâu mà INLA duy trì phổ
rộng hơn baseline, tức tốc độ suy giảm effective rank qua các lớp bị chậm lại.
\end{theorem}
\begin{proof}[Phác thảo]
Kết quả mang tính cục bộ và điều kiện. Baseline vận hành trên feature map hạn chế nên khi xếp chồng
sâu, các vector key/query dễ hội tụ vào vài hướng chủ đạo. INLA thêm lifting phi tuyến trước
aggregation, tăng số basis truy cập được trong không gian trung gian. Nếu regularization đủ ngăn
bùng nổ biên độ và $r > d_k$ nhưng chưa gây nhiễu mạnh, attention duy trì nhiều thành phần độc lập
hơn, nên effective rank giảm chậm hơn theo độ sâu. Đây là lập luận cấu trúc, được kiểm tra bằng
ablation và đo singular value decay trong thực nghiệm (Mục~\ref{sec:exp}).
\end{proof}

\subsection{Giả thuyết thực nghiệm}
\begin{quote}
\textbf{H1:} Tăng hạng của bộ đệm KV qua không gian mở rộng giúp INLA duy trì phổ singular value dồi
dào hơn và chặn tốc độ suy giảm hạng so với linear attention cổ điển khi độ sâu tăng.
\end{quote}
\begin{quote}
\textbf{H2:} Cơ chế lifting học được làm giảm oversmoothing và tăng attention diversity, thể hiện qua
độ phân tán attention, entropy phổ và độ ổn định gradient.
\end{quote}
\begin{quote}
\textbf{H3:} Lợi ích cấu trúc không chuyển hóa thành cải thiện downstream đồng đều ở mọi setting; tác
động của INLA phụ thuộc mạnh vào depth, sequence length, nhiệm vụ và mức regularization.
\end{quote}

% =========================================================
\section{Phương pháp đề xuất: INLA}
\label{sec:method}

\subsection{Tổng quan}
INLA được xây dựng theo tinh thần inverted bottleneck:
\begin{center}
\textit{compression $\rightarrow$ nonlinearity $\rightarrow$ expansion $\rightarrow$ linear attention aggregation}
\end{center}
Một tầng phi tuyến được đưa vào giữa hai bước tuyến tính có kiểm soát, buộc attention đi qua một
không gian trung gian giàu hơn thay vì hoạt động ngay trên biểu diễn gốc.

\subsection{Ánh xạ lifting}
Cho $\mathbf{X} \in \mathbb{R}^{N \times d}$, định nghĩa
\begin{equation}
\Phi_{\mathrm{INLA}}(\mathbf{X}) = \sigma\!\left(\mathbf{X}\mathbf{W}_{\mathrm{low}} + \mathbf{1}\mathbf{b}_{\mathrm{low}}^T\right)\mathbf{W}_{\mathrm{exp}} + \mathbf{1}\mathbf{b}_{\mathrm{exp}}^T,
\end{equation}
với $\mathbf{W}_{\mathrm{low}} \in \mathbb{R}^{d \times d_k}$ (nén), $\mathbf{W}_{\mathrm{exp}} \in
\mathbb{R}^{d_k \times r}$ (mở rộng), $\sigma$ là GeLU/SiLU, $d_k < d$ và $r > d$. Lifting áp dụng
cho query và key, còn value giữ nguyên:
\begin{equation}
\hat{\mathbf{Q}} = \Phi_{\mathrm{INLA}}(\mathbf{Q}), \quad
\hat{\mathbf{K}} = \Phi_{\mathrm{INLA}}(\mathbf{K}), \quad
\hat{\mathbf{V}} = \mathbf{V}.
\end{equation}

\subsection{Phép tính attention tuyến tính}
Khai thác tính kết hợp của nhân ma trận để tránh ma trận $N\times N$:
\begin{equation}
\mathbf{S} = \hat{\mathbf{K}}^T\mathbf{V} \in \mathbb{R}^{r \times d_v}
\quad\text{(nén ngữ cảnh toàn cục)},\qquad
\mathbf{O} = \mathbf{D}^{-1}\big(\hat{\mathbf{Q}}\mathbf{S}\big)
\quad\text{(truy xuất theo query)},
\end{equation}
với chuẩn hóa $\mathbf{d} = \hat{\mathbf{Q}}\big(\hat{\mathbf{K}}^T\mathbf{1}_N\big) \in \mathbb{R}^N$
và $\mathbf{D} = \operatorname{diag}(\mathbf{d})$. Để mẫu số luôn dương và ổn định lan truyền ngược,
ta áp một kernel feature map không âm (cụ thể $\psi(z)=\mathrm{ELU}(z)+1$) lên $\hat{\mathbf{Q}},
\hat{\mathbf{K}}$ sau lifting.

\subsection{Phân tích độ phức tạp}
Với $\mathbf{X} \in \mathbb{R}^{N \times d}$, chi phí một block INLA là
\begin{equation}
\mathcal{O}_{\mathrm{INLA}} = \mathcal{O}\!\Big(N\cdot\big(2dd_k + 2d_kr + 2rd_v\big)\Big),
\end{equation}
trong đó $d, d_k, r, d_v$ là siêu tham số cố định. Do đó độ phức tạp theo độ dài chuỗi vẫn là
$\mathcal{O}(N)$ --- điều kiện nền tảng để INLA phù hợp ngữ cảnh mobile/edge. Tính chất này được
\emph{xác nhận bằng đo đạc} trong Mục~\ref{sec:exp}.

\subsection{Khác biệt với linear attention cổ điển}
Linear attention cổ điển dùng feature map cố định hoặc ít tham số nên khả năng điều chỉnh cấu trúc
phổ hạn chế. INLA bổ sung một \emph{learned feature map có kiểm soát}: lifting xảy ra \emph{trước}
aggregation và tác động trực tiếp lên không gian mà attention vận hành, nhằm điều chỉnh hình học
biểu diễn chứ không chỉ tăng số tham số.

% =========================================================
\section{Hiện thực trong codebase}
\label{sec:impl}

Tính mới được hiện thực thành module thật trong dự án, ở \texttt{cv\_nets/blocks/inla.py}, và export
qua \texttt{cv\_nets.blocks}. Trích đoạn cốt lõi (đã rút gọn):

\begin{lstlisting}
class INLAAttention(nn.Module):
    def __init__(self, dim, dim_compress=None, dim_expand=None,
                 act="gelu", use_lifting=True, ...):
        dim_compress = dim_compress or max(1, dim // 2)   # d_k < d
        dim_expand   = dim_expand   or (2 * dim)          # r  > d
        self.q_proj = nn.Linear(dim, dim); self.k_proj = nn.Linear(dim, dim)
        self.v_proj = nn.Linear(dim, dim)
        if use_lifting:                  # inverted bottleneck: d -> d_k -> (sigma) -> r
            self.lift = nn.Sequential(nn.Linear(dim, dim_compress),
                                      get_activation(act),
                                      nn.Linear(dim_compress, dim_expand))

    def _phi(self, x):                   # lifting + kernel feature map khong am
        return F.elu(self.lift(x)) + 1.0

    def forward(self, x):                # x: (B, N, d)
        q = self._phi(self.q_proj(x)); k = self._phi(self.k_proj(x)); v = self.v_proj(x)
        context = torch.einsum("bnf,bnd->bfd", k, v)        # S = K^T V  (B, F, d)
        denom   = torch.einsum("bnf,bf->bn", q, k.sum(1)).clamp_min(eps).unsqueeze(-1)
        out     = torch.einsum("bnf,bfd->bnd", q, context) / denom   # D^-1 (Q S)
        return self.out_proj(out)
\end{lstlisting}

Một điểm thiết kế quan trọng phục vụ \emph{ablation công bằng}: cờ \texttt{use\_lifting=False} cho ra
baseline linear attention dùng \emph{cùng} pipeline (cùng chuẩn hóa, cùng kernel feature map) nhưng
không có lifting --- nhờ đó mọi khác biệt đo được quy đúng về tác động của lifting. Module
\texttt{INLATransformerBlock} bọc INLA trong khối pre-norm chuẩn (residual + MLP + DropPath).

Ngoài ra, công cụ chẩn đoán phổ được hiện thực ở \texttt{cv\_nets/utils/spectral.py}:
\texttt{effective\_rank}, \texttt{spectral\_entropy}, \texttt{singular\_value\_decay}. Các hàm này
được kiểm định đúng đắn: với ma trận ngẫu nhiên $64\times 48$, $\mathrm{erank}=41.28$ (gần cận trên
$48$); với ma trận hạng-4 dựng cố ý, $\mathrm{erank}=3.96 \approx 4$ (sai số do nhiễu số học).

% =========================================================
\section{Kết quả thực nghiệm sơ bộ}
\label{sec:exp}

Mọi số liệu dưới đây được \emph{đo trực tiếp} trên mã nguồn dự án (script
\texttt{tools/exp\_rank\_collapse.py} và các probe đi kèm), nhằm kiểm chứng những phát biểu định
lượng nhất trước khi bước vào huấn luyện đầy đủ.

\subsection{Độ phức tạp được xác nhận tuyến tính theo $N$}
Bảng~\ref{tab:complexity} đo độ trễ một block INLA (CPU, 4 luồng, $d=64$, $d_k=32$, $r=128$) khi số
token tăng $32\times$ (từ $196$ lên $6272$). Thời gian trên mỗi token gần như phẳng, độ trễ tổng chỉ
tăng $\sim 20\times$ chứ không phải $\sim 1024\times$ như attention bậc hai --- khẳng định
$\mathcal{O}(N)$.

\begin{table}[H]
\centering
\caption{Độ trễ theo số token (CPU, 4 luồng). Baseline = INLA với \texttt{use\_lifting=False}.}
\label{tab:complexity}
\begin{tabular}{@{}rrrrr@{}}
\toprule
$N$ & base (ms) & INLA (ms) & base ($\mu$s/token) & INLA ($\mu$s/token) \\
\midrule
196  & 0.255 & 0.375 & 1.30 & 1.91 \\
784  & 0.420 & 1.060 & 0.54 & 1.35 \\
1568 & 1.290 & 2.093 & 0.82 & 1.33 \\
3136 & 2.508 & 3.665 & 0.80 & 1.17 \\
6272 & 4.709 & 8.178 & 0.75 & 1.30 \\
\bottomrule
\end{tabular}
\end{table}

\subsection{Chi phí tham số khiêm tốn}
Ở $d=64$, một block baseline có $16\,640$ tham số; INLA ($d_k=32$, $r=128$) có $22\,944$ tham số ---
tăng thêm $6\,304$ ($+37.9\%$ cho riêng module attention). Trong toàn backbone (phần lớn tham số nằm
ở convolution), tỷ lệ tăng tổng thể nhỏ hơn nhiều.

\subsection{Thí nghiệm thăm dò rank collapse}
Theo tinh thần Dong et al.\ \cite{dong2021rankcollapse}, ta xếp chồng $L$ khối và đo effective rank
của biểu diễn token ($N=196$, $d=64$, trung bình 8 seed) theo độ sâu, ở hai chế độ.

\paragraph{Chế độ attention thuần (cô lập aggregation, không residual/MLP).}
Cả baseline lẫn INLA đều suy giảm hạng rất nhanh: effective rank rơi từ $\sim 30$ ở lớp 1 xuống
$\approx 1$ quanh lớp 8--9. Điều này tái lập đúng kết quả lý thuyết ``pure attention loses rank'' và
cho thấy: khi cô lập, \emph{cơ chế aggregation tự nó} đẩy biểu diễn về hạng 1 bất kể loại attention.

\paragraph{Chế độ khối transformer thực tế (residual + MLP).}
Bảng~\ref{tab:rank} cho thấy ở khởi tạo ngẫu nhiên, \emph{dòng residual chi phối} và giữ effective
rank cao ($\sim 31\to 34$, thậm chí tăng nhẹ theo độ sâu); INLA và baseline gần như trùng nhau. Đây
là một phát hiện trung thực và nhất quán với lý thuyết: rank collapse là hiện tượng của \emph{pure}
attention, bị residual làm dịu; do đó lợi ích biểu diễn của INLA là \emph{hiệu ứng của chế độ đã
huấn luyện} (khi trọng số lifting được tối ưu), chứ không phải hiệu ứng ở khởi tạo.

\begin{table}[H]
\centering
\caption{Effective rank theo độ sâu trong khối transformer thực tế (residual+MLP), trung bình 8 seed.
Cận trên lý thuyết $\min(N,d)=64$.}
\label{tab:rank}
\begin{tabular}{@{}rrrr@{}}
\toprule
Độ sâu & erank (baseline) & erank (INLA) & spectral entropy (INLA) \\
\midrule
1  & 30.79 & 30.77 & 0.824 \\
4  & 31.98 & 31.98 & 0.833 \\
8  & 32.91 & 32.91 & 0.840 \\
12 & 33.55 & 33.57 & 0.845 \\
16 & 34.09 & 34.11 & 0.849 \\
\bottomrule
\end{tabular}
\end{table}

\subsection{Diễn giải}
Kết quả ủng hộ giả thuyết H3 một cách rõ ràng: tác động của INLA \emph{có điều kiện}. Hạ tầng chẩn
đoán (effective rank, spectral entropy, singular value decay) và thiết kế ablation
(\texttt{use\_lifting} on/off) đã sẵn sàng để kiểm tra H1, H2 \emph{trong quá trình huấn luyện thật},
nơi trọng số lifting được tối ưu theo dữ liệu --- đó là bước tiếp theo tự nhiên của dự án.

% =========================================================
\section{Tích hợp INLA vào backbone nhẹ}

\subsection{Vị trí}
INLA phù hợp nhất khi thay/bổ sung cho các stage \emph{trung gian--sâu} của backbone, nơi số token
đã giảm (sau vài lần hạ mẫu) và receptive field đủ rộng để attention toàn cục có ý nghĩa. Stem và các
stage đầu nên giữ thiên hướng convolution để tiết kiệm chi phí; head cuối dùng pooling/projection
nhẹ.

\subsection{Các biến thiết kế cần kiểm soát}
$d_k$ (chiều nén), $r$ (chiều mở rộng), loại activation $\sigma$, vị trí normalization, số head, và
regularization trên ma trận chiếu/đầu ra expansion. Cấu hình mặc định đề xuất: $d_k$ đủ nhỏ để rẻ
nhưng không quá hẹp; $r > d$ ở mức vừa phải để tránh khuếch đại nhiễu; GeLU/SiLU; normalization ngay
sau aggregation; regularization nhẹ trên lifting để kiểm soát trade-off rank--efficiency.

% =========================================================
\section{Kế hoạch thực nghiệm đầy đủ}

\subsection{Baseline cho ablation công bằng}
Cùng backbone, cùng số stage/hidden size, gần tương đương FLOPs, gồm: (1) linear attention tiêu
chuẩn; (2) kernel feature map cố định; (3) lifting tuyến tính (bỏ phi tuyến); (4) phi tuyến nhưng
không expansion; (5) expansion nhưng không regularization. Trong codebase, các biến (1) và (toàn bộ
INLA) đã có sẵn qua cờ \texttt{use\_lifting}; các biến còn lại là chỉnh sửa nhỏ trên \texttt{lift}.

\subsection{Chỉ số đánh giá}
\begin{table}[H]
\centering
\caption{Khung đánh giá đề xuất}
\label{tab:metrics}
\begin{tabular}{@{}p{4cm}p{10cm}@{}}
\toprule
\textbf{Nhóm} & \textbf{Chỉ số} \\
\midrule
Cấu trúc biểu diễn & Effective rank, spectral entropy, singular value decay \\
Hành vi attention & Diversity, oversmoothing, gradient norm/variance, loss smoothness \\
Hiệu năng hệ thống & Latency, throughput, peak memory, FLOPs, số tham số \\
Hiệu năng tác vụ & IoU/Top-1/F1 tùy bài toán \\
\bottomrule
\end{tabular}
\end{table}

% =========================================================
\section{Phân tích failure modes}

\begin{itemize}
    \item \textbf{$r$ quá lớn:} lifting khuếch đại nhiễu thay vì trích đặc trưng; hạng đại số tăng
    nhưng tín hiệu hữu ích không tăng tương ứng $\Rightarrow$ overfit.
    \item \textbf{Backbone quá nông:} rank collapse chưa đủ nghiêm trọng để INLA tạo lợi ích; lifting
    chỉ thêm chi phí. (Phù hợp với kết quả ở Mục~\ref{sec:exp}: ở khởi tạo/độ sâu vừa, khác biệt nhỏ.)
    \item \textbf{Task cục bộ:} nếu nhiệm vụ chủ yếu dựa đặc trưng cục bộ, basis attention giàu hơn
    không tương xứng lợi ích; backbone thuần conv có thể cạnh tranh tốt hơn.
    \item \textbf{Regularization sai mức:} quá mạnh thì lifting mất tác dụng; quá yếu thì mất kiểm
    soát trade-off rank--efficiency.
\end{itemize}

Ngay cả khi INLA không cải thiện mạnh downstream, dự án vẫn có giá trị nếu chỉ ra được: cấu trúc phổ
thay đổi đúng dự báo, miền hoạt động hiệu quả của nonlinear lifting, và hướng dẫn thiết kế attention
nhẹ.

% =========================================================
\section{Lộ trình thực hiện}
\begin{enumerate}[leftmargin=*]
    \item Rà soát tài liệu, chốt baseline và phạm vi tác vụ, hoàn thiện câu hỏi/giả thuyết.
    \item Dựng backbone cơ sở (MV2Block + linear attention), cố định cấu hình, đảm bảo train ổn định.
    \item \textbf{Tích hợp INLA} (đã hoàn thành ở mức module): pipeline compression--nonlinearity--expansion,
    khảo sát activation, $d_k$, $r$, normalization. \emph{(Hoàn tất phần hiện thực + kiểm thử.)}
    \item Ablation + đo cấu trúc (effective rank, spectral entropy, singular value decay) theo depth
    và sequence length, nhiều seed. \emph{(Hạ tầng đã sẵn sàng.)}
    \item Đối chiếu thực nghiệm với các mệnh đề; phân tích failure modes bằng dữ liệu.
    \item Đánh giá hiệu năng thực tế (IoU/latency/memory/FLOPs) và trade-off.
    \item Tổng hợp, hoàn thiện báo cáo, điều chỉnh limitation để tránh overclaim.
\end{enumerate}

% =========================================================
\section{Kết luận}

Báo cáo định vị INLA không phải sự thay thế đại trà cho mọi attention, mà là một \emph{can thiệp có
kiểm soát} lên cấu trúc phổ của linear attention trong backbone thị giác gọn nhẹ. Chúng tôi đã: (i)
mô tả cơ chế rank collapse và động cơ thiết kế; (ii) xây dựng các định nghĩa, mệnh đề, bổ đề và một
định lý cục bộ giải thích vì sao INLA có thể làm chậm suy giảm hạng trong miền giả thiết rõ ràng;
(iii) \emph{hiện thực INLA thành module thật} trong codebase, kèm công cụ chẩn đoán phổ và thiết kế
ablation công bằng; (iv) báo cáo các số liệu đo thật: độ phức tạp được xác nhận tuyến tính theo số
token, chi phí tham số khiêm tốn, và một thí nghiệm thăm dò rank collapse trung thực.

Phát hiện thực nghiệm sơ bộ nhấn mạnh tính \emph{có điều kiện} của lợi ích (giả thuyết H3): ở khởi
tạo ngẫu nhiên, residual che khuất khác biệt, trong khi cơ chế collapse lộ rõ ở attention thuần. Vì
trọng số lifting là tham số học được, kiểm chứng H1/H2 thuộc về chế độ đã huấn luyện --- và toàn bộ
hạ tầng đo lường cần thiết đã được dựng sẵn. Chính sự minh bạch về điều kiện áp dụng làm cho nghiên
cứu có giá trị học thuật thực sự, và mở ra một cách nhìn về thiết kế attention nhẹ: không chỉ tối ưu
FLOPs, mà còn điều khiển hình học biểu diễn bên trong attention.

% =========================================================
\begin{thebibliography}{99}

\bibitem{vaswani2017attention}
A. Vaswani et al., ``Attention is all you need,'' in \textit{NeurIPS}, 2017.

\bibitem{mobilenetv2}
M. Sandler, A. Howard, M. Zhu, A. Zhmoginov, and L.-C. Chen, ``MobileNetV2: Inverted residuals and linear bottlenecks,'' in \textit{CVPR}, 2018.

\bibitem{mobilevitv1}
S. Mehta and M. Rastegari, ``MobileViT: Light-weight, general-purpose, and mobile-friendly vision transformer,'' in \textit{ICLR}, 2022.

\bibitem{mobilevitv2}
S. Mehta and M. Rastegari, ``Separable self-attention for mobile vision transformers,'' \textit{TMLR}, 2023.

\bibitem{katharopoulos2020transformers}
A. Katharopoulos, A. Vyas, N. Pappas, and F. Fleuret, ``Transformers are RNNs: Fast autoregressive transformers with linear attention,'' in \textit{ICML}, 2020.

\bibitem{shen2020efficient}
Z. Shen, M. Zhang, H. Zhao, S. Yi, and H. Li, ``Efficient attention: Attention with linear complexities,'' in \textit{WACV}, 2021.

\bibitem{dong2021rankcollapse}
Y. Dong, J.-B. Cordonnier, and A. Loukas, ``Attention is not all you need: Pure attention loses rank doubly exponentially with depth,'' in \textit{ICML}, 2021.

\bibitem{bhojanapalli2020low}
S. Bhojanapalli, C. Yun, A. S. Rawat, S. Reddi, and S. Kumar, ``Low-rank bottleneck in multi-head attention models,'' in \textit{ICML}, 2020.

\bibitem{choromanski2021performer}
K. Choromanski et al., ``Rethinking attention with performers,'' in \textit{ICLR}, 2021.

\bibitem{wang2020linformer}
S. Wang, B. Z. Li, M. Khabsa, H. Fang, and H. Ma, ``Linformer: Self-attention with linear complexity,'' \textit{arXiv:2006.04768}, 2020.

\bibitem{xiong2021nystromformer}
Y. Xiong et al., ``Nyströmformer: A Nyström-based algorithm for approximating self-attention,'' in \textit{AAAI}, 2021.

\bibitem{qin2022cosformer}
Z. Qin et al., ``cosFormer: Rethinking softmax in attention,'' in \textit{ICLR}, 2022.

\end{thebibliography}

\end{document}
